\begin{definition}[Heat semigroup Besov space]
	Let $P_t = e^{-tL}$ be the Markov semigroup generated by $L$.
	For $s>0$ and $1\le p,q\le\infty$, define
	\[
		\|f\|_{B^s_{p,q}(L)}
		=
		\|f\|_{L^p(\mu)}
		+
		\left(
		\int_0^1
		\left(
		t^{-s/2}
		\| (I-P_t) f \|_{L^p(\mu)}
		\right)^q
		\frac{dt}{t}
		\right)^{1/q}.
	\]
	Then $B^s_{p,q}(L)$ consists of all $f$ for which this norm is finite.
\end{definition}

\begin{definition}[Spectral wavelets associated to $L$]
	Let $\varphi_0 \in C^\infty_c([0,\infty))$ satisfy
	$\varphi_0(\lambda) = 1$ for $\lambda \in [0,1]$
	and $\operatorname{supp} \varphi_0 \subset [0,2]$.
	Define the dyadic wavelet kernel
	\[
		\psi(\lambda) = \varphi_0(\lambda) - \varphi_0(2\lambda),
		\qquad
		\psi_j(\lambda) = \psi(2^{-j}\lambda), \quad j \ge 1,
	\]
	so that $\{\varphi_0, \psi_j\}_{j \ge 1}$ forms a partition of unity:
	\[
		\varphi_0(\lambda) + \sum_{j=1}^{\infty} \psi_j(\lambda) = 1
		\qquad \text{for all } \lambda \ge 0.
	\]
	The \emph{wavelet blocks} are the operators
	\[
		\Psi_0 = \varphi_0(L),
		\qquad
		\Psi_j = \psi_j(L) = \psi(2^{-j}L), \quad j \ge 1,
	\]
	defined via the spectral calculus of $L$, yielding the
	Littlewood--Paley decomposition
	\[
		f = \Psi_0 f + \sum_{j=1}^{\infty} \Psi_j f.
	\]
	If $L$ has discrete spectrum $0 \le \lambda_1 \le \lambda_2 \le \cdots$
	with eigenfunctions $\{e_k\}$, then
	\[
		\Psi_j f
		= \sum_{k} \psi_j(\lambda_k)\, \langle f, e_k \rangle_{L^2(\mu)}\, e_k.
	\]
\end{definition}
















% \begin{definition}[Spectral Sobolev class]
% 	For smoothness $s > 0$ and radius $R > 0$, the \emph{spectral Sobolev class}
% 	of order $s$ is
% 	\[
% 		\mathcal{H}^s_\mu(R)
% 		:=
% 		\Bigl\{
% 		f \in L^2(\mu)
% 		:
% 		\| f \|_{\mathcal{H}^s}^2
% 		:=
% 		\sum_{j \geq 1} \lambda_j^{s}\, \hat{f}_j^{\,2}
% 		\leq R^2
% 		\Bigr\},
% 		\qquad
% 		\hat{f}_j := \langle f, \phi_j \rangle_{L^2(\mu)}.
% 	\]
% 	The condition $f \in \mathcal{H}^s_\mu(R)$ requires the spectral coefficients
% 	to decay as $|\hat{f}_j| \leq R\, \lambda_j^{-s/2} \asymp R\, j^{-s/d_{\mathrm{eff}}}$:
% 	higher smoothness $s$ enforces faster decay and hence greater regularity
% 	with respect to the geometry encoded by $\mathcal{E}$.
% \end{definition}


% \begin{definition}[Heat semigroup Besov space]
% 	Let $P_t = e^{-tL}$ be the Markov semigroup generated by $L$.
% 	For $s>0$ and $1\le p,q\le\infty$, define
% 	\[
% 		\|f\|_{B^s_{p,q}(L)}
% 		=
% 		\|f\|_{L^p(\mu)}
% 		+
% 		\left(
% 		\int_0^1
% 		\left(
% 		t^{-s/2}
% 		\| (I-P_t) f \|_{L^p(\mu)}
% 		\right)^q
% 		\frac{dt}{t}
% 		\right)^{1/q}.
% 	\]
% 	Then $B^s_{p,q}(L)$ consists of all $f$ for which this norm is finite.
% \end{definition}

% \begin{definition}[Spectral wavelets associated to $L$]
% 	Let $\varphi_0 \in C^\infty_c([0,\infty))$ satisfy
% 	$\varphi_0(\lambda) = 1$ for $\lambda \in [0,1]$
% 	and $\operatorname{supp} \varphi_0 \subset [0,2]$.
% 	Define the dyadic wavelet kernel
% 	\[
% 		\psi(\lambda) = \varphi_0(\lambda) - \varphi_0(2\lambda),
% 		\qquad
% 		\psi_j(\lambda) = \psi(2^{-j}\lambda), \quad j \ge 1,
% 	\]
% 	so that $\{\varphi_0, \psi_j\}_{j \ge 1}$ forms a partition of unity:
% 	\[
% 		\varphi_0(\lambda) + \sum_{j=1}^{\infty} \psi_j(\lambda) = 1
% 		\qquad \text{for all } \lambda \ge 0.
% 	\]
% 	The \emph{wavelet blocks} are the operators
% 	\[
% 		\Psi_0 = \varphi_0(L),
% 		\qquad
% 		\Psi_j = \psi_j(L) = \psi(2^{-j}L), \quad j \ge 1,
% 	\]
% 	defined via the spectral calculus of $L$, yielding the
% 	Littlewood--Paley decomposition
% 	\[
% 		f = \Psi_0 f + \sum_{j=1}^{\infty} \Psi_j f.
% 	\]
% 	If $L$ has discrete spectrum $0 \le \lambda_1 \le \lambda_2 \le \cdots$
% 	with eigenfunctions $\{e_k\}$, then
% 	\[
% 		\Psi_j f
% 		= \sum_{k} \psi_j(\lambda_k)\, \langle f, e_k \rangle_{L^2(\mu)}\, e_k.
% 	\]
% \end{definition}



% \begin{definition}[Spectral Besov space associated with a Dirichlet form]
% 	Let $(X,\mathcal{F},\mu,\mathcal{E})$ be a measure space endowed with a Dirichlet form,
% 	and let $L$ be the associated non-negative self-adjoint operator on $L^2(\mu)$
% 	with spectral resolution
% 	\[
% 		L\phi_j = \lambda_j \phi_j,
% 		\qquad
% 		0=\lambda_0 \le \lambda_1 \le \lambda_2 \le \cdots.
% 	\]

% 	For $s \in \mathbb{R}$ and $1 \le p,q \le \infty$, define dyadic spectral projectors
% 	\[
% 		\Delta_k f
% 		:=
% 		\sum_{2^{2k} \le 1+\lambda_j < 2^{2(k+1)}}
% 		\langle f,\phi_j\rangle_{L^2(\mu)} \phi_j.
% 	\]

% 	The Besov space $B^s_{p,q}(L)$ consists of all
% 	$f \in L^p(\mu)$ such that
% 	\[
% 		\|f\|_{B^s_{p,q}(L)}
% 		=
% 		\begin{cases}
% 			\left(
% 			\displaystyle \sum_{k=0}^\infty
% 			2^{ksq}
% 			\|\Delta_k f\|_{L^p(\mu)}^q
% 			\right)^{1/q},
% 			 & 1 \le q < \infty, \\[1em]
% 			\displaystyle \sup_{k \ge 0}
% 			2^{ks}
% 			\|\Delta_k f\|_{L^p(\mu)},
% 			 & q = \infty,
% 		\end{cases}
% 	\]
% 	is finite.
% \end{definition}



% Our learning guarantees hold if the effective dimension is small. This is well-know as bias-variance tradeoff.






% \begin{proposition}[Hypothesis 1:]
% 	Given observations
% \end{proposition}

% \begin{proposition}[Laplacian and spectrum]
% Given a learning space $(\cX, \mu, \cE)$, there exists a unique non-negative self-adjoint operator $\cL$ on $L^2(\cX, \mu)$ such that
% \[
% 	\cE(f, g) = \langle f, \cL g \rangle_{L^2(\mu)}, \qquad \forall\, f \in \cD(\cE),\; g \in \cD(\cL).
% \]
% The operator $\cL$ admits a spectral decomposition $\cL \phi_j = \lambda_j \phi_j$, with
% \[
% 	0 = \lambda_0 \leq \lambda_1 \leq \lambda_2 \leq \cdots, \qquad \langle \phi_j, \phi_k \rangle_{L^2(\mu)} = \delta_{jk}.
% \]
% \end{proposition}


% Equivalently, $\pi_\lambda$ is absolutely continuous with respect to $\nu$, with Radon--Nikodym derivative
% \[
% 	\frac{d\pi_\lambda}{d\nu}(f)
% 	=
% 	\frac{1}{Z_\lambda} e^{-\lambda R(f)}.
% \]


% \begin{remark}
% 	\emph{Given observations $y$ with likelihood $p(y \mid f)$, the posterior distribution induced by the Gibbs prior is $\pi_\lambda(\mathrm df \mid y) \propto	p(y \mid f)  e^{-\lambda R(f)} \nu(\mathrm df)$. The maximum a posteriori (MAP) estimator is therefore given by
% 		\begin{equation}
% 			\hat f_{\mathrm{MAP}}
% 			=
% 			\arg\min_{f \in X}
% 			\Big(
% 			- \log p(y \mid f) + \lambda R(f)
% 			\Big),
% 		\end{equation}
% 		which coincides with regularized empirical risk minimization.}
% \end{remark}


% \section{Effective Dimension Hypothesis}

% Let us formulate the extension of the manifold hypothesis to this case. \emile{more generally, need to define energy functions, why they're relevant to this problem, and what characteristics do we want our manifold hypothesis to have?}


% \begin{assumption}[Effective Dimension Hypothesis] \label{hyp:eff_dim}
% 	All the observed data $x_i$ come from a learning space $(X, \mathcal{F}, \mu, \mathcal{E})$ with small effective dimension even if the ambient space  $(X, \mathcal{F}, \nu_0, \mathcal{E})$ has large effective dimension where $\nu_0$ is the base distribution on $\mathcal{X}$:
% 	\[
% 		d_{\mathrm{eff}}(X, \mathcal{F}, \mu, \mathcal{E}) \ll d_{\mathrm{eff}}(X, \mathcal{F}, \nu_0, \mathcal{E})
% 	\]
% 	% and
% 	% \[
% 	% 	D_{KL}(\mu \Vert \nu_0) = \int \log \frac{d\mu}{d\nu_0} d\mu < \infty
% 	% \]
% \end{assumption}

% This assumption ensures two things: existence of the efficient coordinate system in the space $\mathcal{X}$ and efficient learning rate for functions in Besov spaces. Moreover, this assumption of the underlying learning space is strictly stronger than the manifold hypothesis and unifies multiple learning problems, e.g. classical kernel methods, flow matching/diffusion models and transformers.


% \subsection{Generalization Error}


% \begin{proposition}[Bias-variance tradeoff]
% 	Let $f^* \in \mathcal{H}^s_\mu(R)$ and suppose we observe
% 	$y_i = f^*(x_i) + \varepsilon_i$, $i = 1, \ldots, n$,
% 	with $x_i \sim \mu$ i.i.d.\ and $\varepsilon_i$ i.i.d.\
% 	sub-Gaussian with parameter $\sigma^2$.
% 	Consider the truncated spectral estimator
% 	\[
% 		\hat{f}_k(x)
% 		=
% 		\sum_{j=1}^{k} \hat{f}_j^{(n)}\, \phi_j(x),
% 		\qquad
% 		\hat{f}_j^{(n)}
% 		=
% 		\frac{1}{n} \sum_{i=1}^{n} y_i\, \phi_j(x_i).
% 	\]
% 	Then the expected risk decomposes as
% 	\[
% 		\mathbb{E}\bigl[\|\hat{f}_k - f^*\|_{L^2(\mu)}^2\bigr]
% 		=
% 		\underbrace{
% 			\sum_{j > k} (\hat{f}^*_j)^2
% 		}_{\mathrm{bias}^2(k)}
% 		\;+\;
% 		\underbrace{
% 			\frac{\sigma^2\, k}{n}
% 		}_{\mathrm{variance}(k)}.
% 	\]
% 	Under Effective Dimension Hypothesis (\Cref{hyp:eff_dim}) ($\lambda_j \asymp j^{2/d_{\mathrm{eff}}}$)
% 	and the Sobolev constraint $\sum_j \lambda_j^s (\hat{f}^*_j)^2 \leq R^2$:

% 	\medskip
% 	\noindent\textbf{Bias bound.}
% 	The tail energy is controlled by the smoothness and dimension:
% 	\[
% 		\mathrm{bias}^2(k)
% 		=
% 		\sum_{j > k} (\hat{f}^*_j)^2
% 		\leq
% 		R^2\, \lambda_{k+1}^{-s}
% 		\asymp
% 		R^2\, k^{-2s/d_{\mathrm{eff}}}.
% 	\]

% 	\noindent\textbf{Variance bound.}
% 	\[
% 		\mathrm{variance}(k) = \frac{\sigma^2\, k}{n}.
% 	\]

% 	\noindent\textbf{Optimal truncation.}
% 	Balancing bias and variance:
% 	\[
% 		k^{-2s/d_{\mathrm{eff}}} \asymp \frac{k}{n}
% 		\qquad\Longrightarrow\qquad
% 		k^* \asymp n^{d_{\mathrm{eff}}/(2s + d_{\mathrm{eff}})}.
% 	\]

% 	\noindent\textbf{Minimax rate.}
% 	Substituting $k^*$ into either term:
% 	\[
% 		\boxed{
% 		\mathbb{E}\bigl[\|\hat{f}_{k^*} - f^*\|_{L^2(\mu)}^2\bigr]
% 		\;\lesssim\;
% 		n^{-2s/(2s + d_{\mathrm{eff}})}.
% 		}
% 	\]
% 	This rate is minimax optimal over $\mathcal{H}^s_\mu(R)$.
% \end{proposition}

% \begin{remark}[Interpretation]
% 	The three quantities $s$, $d_{\mathrm{eff}}$, and $n$ interact as follows:
% 	\begin{itemize}
% 		\item \textbf{Smoothness $s$}:
% 		      higher $s$ means faster eigenvalue decay of the
% 		      coefficients $\hat{f}^*_j$, so the bias
% 		      decreases faster with $k$ and fewer modes suffice.
% 		\item \textbf{Effective dimension $d_{\mathrm{eff}}$}:
% 		      larger $d_{\mathrm{eff}}$ means eigenvalues
% 		      $\lambda_j \asymp j^{2/d_{\mathrm{eff}}}$ grow
% 		      more slowly, so more modes carry significant energy
% 		      and the variance cost of including them is higher.
% 		      This is the curse of dimensionality in $d_{\mathrm{eff}}$,
% 		      not in the ambient dimension.
% 		\item \textbf{Optimal number of modes $k^*$}:
% 		      this is the resolution at which you can reliably
% 		      estimate the function given $n$ samples.
% 		      Modes $j \leq k^*$ are estimated;
% 		      modes $j > k^*$ are discarded as noise.
% 	\end{itemize}
% \end{remark}


% \begin{theorem}[Minimax rate over Besov balls \cite{donohoDensityEstimationWavelet1996}]

% 	\todo{Replace $\Real^d$ by $\mathcal{X}$}
% 	Let $f \in \mathcal{B}^s_{p,q}(R) \subset L^r(\mathbb{R}^d)$ with
% 	$s > d(1/p - 1/r)_+$. Then the minimax $L^r$ estimation rate from
% 	$n$ noisy samples is
% 	\[
% 		\inf_{\hat{f}_n}\, \sup_{f \in \mathcal{B}^s_{p,q}(R)}
% 		\mathbb{E}\bigl[\|\hat{f}_n - f\|_{L^r}\bigr]
% 		\asymp
% 		\begin{cases}
% 			n^{-\frac{s}{2s+d}},
% 			 & \text{if } p \ge r \quad \text{(dense zone)}, \\[6pt]
% 			\displaystyle\left(\frac{\log n}{n}\right)^{\!\frac{s-d(1/p-1/r)}{2(s-d(1/p-1/r))+d}},
% 			 & \text{if } p < r \quad \text{(sparse zone)}.
% 		\end{cases}
% 	\]
% 	The dense rate is achieved by linear methods (projection, kernel smoothing).
% 	The sparse rate requires nonlinear wavelet thresholding; linear estimators
% 	are suboptimal by a polynomial factor in $\log n$.
% \end{theorem}




% \subsection{Compression}


% \begin{theorem}
% 	Lower bound on the number of bits to represent a function $f$ up to the accuracy $\epsilon$ is given by
% 	\begin{align*}
% 		bits(f, \epsilon) \geq \ldots
% 	\end{align*}

% \end{theorem}